% from ega4-10-fr.tex
\oldpage[IV-3]{97}
donc aussi dans $X_0$). En vertu de \emph{c)}, il y a pour tout $\alpha$ un
seul ouvert $U_\alpha$ de $X$ tel que $X_0\cap U_\alpha=V_\alpha$, et en
vertu de \emph{c$''$)} un seul ensemble quasi-constructible $Z_\alpha$ dans
$X$ tel que $Z_0\cap V_\alpha=X_0\cap Z_\alpha$. Si $\alpha$ et $\beta$ sont
deux indices quelconques, on a
\[
Z_\beta\cap U_\alpha\cap X_0
=U_\alpha\cap Z_0\cap V_\beta
=V_\alpha\cap Z_0\cap V_\beta
=Z_\alpha\cap X_0\cap V_\beta
\subset Z_\alpha\cap X_0~;
\]
comme $Z_\beta\cap U_\alpha$ et $Z_\alpha$ sont quasi-constructibles dans
$X$, il résulte de \emph{c$''$)} que
$Z_\beta\cap U_\alpha\subset Z_\alpha$. Si l'on pose
$Z=\bigcup_\alpha Z_\alpha$, on a donc $Z\cap U_\alpha=Z_\alpha$ pour tout
$\alpha$~; d'ailleurs, comme les $V_\alpha$ recouvrent $X_0$, il résulte de
\emph{c)} que les $U_\alpha$ recouvrent $X$, et l'on voit donc que $Z$ est
localement quasi-constructible dans $X$ et $Z_0=Z\cap X_0$.

\begin{definition}[10.1.3]
\label{IV.10.1.3-fr}
Lorsqu'une partie $X_0$ d'un espace topologique vérifie les conditions
équivalentes de (10.1.2), on dit que $X_0$ est \emph{très dense} dans $X$.
\end{definition}

On a déjà vu au cours de la démonstration de (10.1.2) que $X_0$ est alors
\emph{dense} dans $X$.

\begin{corollary}[10.1.4]
\label{IV.10.1.4-fr}
Si $X_0$ est très dense dans $X$ et $U$ une partie ouverte de $X$,
$U\cap X_0$ est très dense dans $U$. Inversement, si $(U_\alpha)$ est un
recouvrement ouvert de $X$ tel que $U_\alpha\cap X_0$ soit très dense dans
$U_\alpha$ pour tout $\alpha$, alors $X_0$ est très dense dans $X$.
\end{corollary}

Comme toute partie localement fermée dans $U$ est localement fermée dans $X$,
la première assertion résulte du critère \emph{a)} de (10.1.2)~; il en est de
même de la seconde, car si $Z\neq\varnothing$ est localement fermée dans $X$,
$Z\cap U_\alpha$ est localement fermée dans $U_\alpha$ pour tout $\alpha$,
et $Z\cap U_\alpha\neq\varnothing$ pour un $\alpha$ au moins.

\subsection{Quasi-homéomorphismes}
\label{subsection:IV.10.2-fr}

\begin{proposition}[10.2.1]
\EGAReaderTarget{ega4.c04.prop.10-2-1}{IV.10.2.1-fr}
Soient $X_0$, $X$ deux espaces topologiques, $f:X_0\to X$ une application
continue. Les conditions suivantes sont équivalentes~:
\begin{enumerate}[label=\emph{\alph*)}]
  \item L'application $U\mapsto f^{-1}(U)$ de $\mathfrak O(X)$ dans
  $\mathfrak O(X_0)$ est bijective.
  \item[\emph{a$'$)}] L'application $F\mapsto f^{-1}(F)$ de $\mathfrak F(X)$
  dans $\mathfrak F(X_0)$ est bijective.
  \item La topologie de $X_0$ est image réciproque par $f$ de celle de $X$, et
  $f(X_0)$ est très dense (10.1.3) dans $X$.
  \item Le foncteur $\mathcal F\mapsto f^*(\mathcal F)$ de la catégorie des
  faisceaux d'ensembles (resp. des faisceaux de groupes abéliens) sur $X$ dans
  la catégorie des faisceaux d'ensembles (resp. des faisceaux de groupes
  abéliens) sur $X_0$ est une équivalence de catégories.
\end{enumerate}
\end{proposition}

Il est clair que \emph{a)} et \emph{a$'$)} sont équivalentes, et \emph{a)}
implique que la topologie de $X_0$ est image réciproque par $f$ de celle de
$X$. D'autre part, si $f(X_0)$ n'est pas très dense dans $X$, il y a deux
ouverts distincts $U_1$, $U_2$ de $X$ tels que
$U_1\cap f(X_0)=U_2\cap f(X_0)$, et par suite
$f^{-1}(U_1)=f^{-1}(U_2)$, ce qui achève de montrer que \emph{a)} entraîne
\emph{b)}. Inversement, la condition \emph{b)} implique que les applications
$U\mapsto U\cap f(X_0)$ de $\mathfrak O(X)$ dans $\mathfrak O(f(X_0))$ et
$V\mapsto f^{-1}(V)$ de $\mathfrak O(f(X_0))$ dans $\mathfrak O(X_0)$ sont
bijectives, donc il en est de même de leur composée $U\mapsto f^{-1}(U)$.

Pour voir que \emph{a)} entraîne \emph{c)}, il suffit d'appliquer la définition
de $f^*(\mathcal F)$ ($\mathbf 0_{\mathrm I}$, 3.7.1) et les axiomes des
faisceaux~: \emph{a)} entraîne que pour tout ouvert $U$ de $X$, l'application
canonique
$\Gamma(U,\mathcal F)\to\Gamma(f^{-1}(U),f^*(\mathcal F))$ est une bijection
fonctorielle en $\mathcal F$, d'où \emph{c)}. Il reste à montrer que
\emph{c)} entraîne \emph{b)}.

